\documentclass[12pt]{article}

\usepackage{amsmath, amssymb, bm, booktabs, geometry, setspace}
\geometry{margin=1in}
\doublespacing

\title{Variance Component Modelling of Nuclear Family Data:\\
From Trio Designs to Full ACDE Estimation in R}
\author{}
\date{}

\begin{document}
\maketitle
\tableofcontents

%====================================================
\section{Introduction}

Biometrical genetics aims to partition phenotypic variation into
genetic and environmental components:

\begin{itemize}
\item Additive genetic effects ($A$)
\item Dominance genetic effects ($D$)
\item Shared environmental effects ($C$)
\item Unique environmental effects ($E$)
\end{itemize}

Twin studies traditionally estimate ACE or ADE models.  
Modern population cohorts, however, frequently contain nuclear family
data. This paper presents a unified framework showing:

\begin{itemize}
\item what can be estimated from trio data,
\item why siblings solve identifiability,
\item how to estimate AE/ACE/ADE/ACDE using linear mixed models in R.
\end{itemize}

%====================================================
\section{General Variance Decomposition}

For phenotype $y$:

\[
y = \mu + a + d + c + e
\]

\[
Var(y) = \sigma_A^2 + \sigma_D^2 + \sigma_C^2 + \sigma_E^2
\]

%====================================================
\section{General Pedigree Covariance}

For relatives $i,k$:

\[
\text{Cov}(y_i,y_k)
= 2\phi_{ik}\sigma_A^2
+ \Delta_{ik}\sigma_D^2
+ \gamma_{ik}\sigma_C^2
\]

\begin{itemize}
\item $\phi_{ik}$ : kinship coefficient
\item $\Delta_{ik}$ : probability of sharing two alleles IBD
\item $\gamma_{ik}$ : shared family environment indicator
\end{itemize}

%====================================================
\section{Nuclear Family Trio Design}

Consider mother--father--child trios.

\subsection{Expected Covariances}

\begin{center}
\begin{tabular}{lccc}
\toprule
Pair & $2\phi$ & $\Delta$ & $\gamma$ \\ \midrule
Mother--Father & 0 & 0 & 1 \\
Mother--Child  & $\tfrac12$ & 0 & 1 \\
Father--Child  & $\tfrac12$ & 0 & 1 \\
\bottomrule
\end{tabular}
\end{center}

\subsection{Key Consequence}

Dominance deviations are not transmitted:

\[
\Delta_{ik}=0 \quad \text{for all trio pairs}
\]

Thus dominance contributes only to individual variance:

\[
Var(y)=\sigma_A^2+\sigma_D^2+\sigma_E^2
\]

\subsection{Identifiable Family Resemblance}

Trio data contain only one familial covariance:

\[
\text{Cov}_{family} = \frac12\sigma_A^2 + \sigma_F^2
\]

\[
\sigma_F^2 =
\begin{cases}
\sigma_C^2 & \text{ACE interpretation} \\
\text{residual familial resemblance} & \text{ADE interpretation}
\end{cases}
\]

Hence:

\[
\textbf{Trio data allow AE, ACE or ADE --- but not ACDE.}
\]

%====================================================
\section{Why Adding Siblings Solves Identifiability}

Sibling covariance:

\[
\text{Cov}(\text{sibling}_1,\text{sibling}_2)
= \tfrac12 A + C + \tfrac14 D
\]

\begin{center}
\begin{tabular}{lc}
\toprule
Design & Identifiable components \\ \midrule
Trio & AE, ACE or ADE \\
Family with siblings & ACDE \\
\bottomrule
\end{tabular}
\end{center}

%====================================================
\section{ACDE Covariance Matrix}

For parents ($P_f,P_m$) and siblings ($S_1,S_2$):

\[
\Sigma =
\begin{pmatrix}
V_P & V_C & \tfrac12 V_A + V_C & \tfrac12 V_A + V_C \\
V_C & V_P & \tfrac12 V_A + V_C & \tfrac12 V_A + V_C \\
\tfrac12 V_A + V_C & \tfrac12 V_A + V_C & V_P &
\tfrac12 V_A + V_C + \tfrac14 V_D \\
\tfrac12 V_A + V_C & \tfrac12 V_A + V_C &
\tfrac12 V_A + V_C + \tfrac14 V_D & V_P
\end{pmatrix}
\]

%====================================================
\section{Likelihood Formulation}

\[
\mathbf{y}_i \sim \mathcal{N}(\mathbf{0}, \Sigma)
\]

\[
\log f(\mathbf{y}_i) =
-\frac12\left[
\log|\Sigma|
+ \mathbf{y}_i^\top \Sigma^{-1} \mathbf{y}_i
+ k\log(2\pi)
\right]
\]

Variance parameterisation:

\[
V_A = e^{\theta_A}, \quad
V_C = e^{\theta_C}, \quad
V_D = e^{\theta_D}, \quad
V_E = e^{\theta_E}
\]

%====================================================
\section{Linear Mixed Model Representation}

Stacked observations:

\[
\mathbf{y} = \mathbf{X}\beta + \mathbf{Z}_A \mathbf{a}
+ \mathbf{Z}_F \mathbf{f} + \boldsymbol{\varepsilon}
\]

\begin{align}
\mathbf{a} &\sim N(0, \sigma_A^2 \mathbf{A}) \\
\mathbf{f} &\sim N(0, \sigma_F^2 \mathbf{I}) \\
\boldsymbol{\varepsilon} &\sim N(0, \sigma_E^2 \mathbf{I})
\end{align}

%====================================================
\section{Data Restructuring in R}

\begin{verbatim}
library(nlme)

longdata <- reshape(
  widelong,
  varying=c("bmi1","bmi2","bmi3"),
  v.names="bmi",
  timevar="member",
  times=c("mother","father","child"),
  idvar="famid",
  direction="long"
)

longdata$Acoef <- ifelse(longdata$member=="child",1,0.5)
longdata$Ccoef <- 1
longdata$Dcoef <- 1
\end{verbatim}

%====================================================
\section{AE Model}

\[
Var(y)=\sigma_A^2+\sigma_E^2
\]

\begin{verbatim}
model_AE <- lme(
  bmi~1,
  random=list(
    famid=pdDiag(~Acoef-1),
    indid=pdDiag(~Dcoef-1)
  ),
  data=longdata, method="REML")
\end{verbatim}

%====================================================
\section{ACE Model}

\[
Var(y)=\sigma_A^2+\sigma_C^2+\sigma_E^2
\]

\begin{verbatim}
model_ACE <- lme(
  bmi~1,
  random=list(
    famid=pdDiag(~Acoef+Ccoef-1),
    indid=pdDiag(~Dcoef-1)
  ),
  data=longdata, method="REML")
\end{verbatim}

%====================================================
\section{ADE Model}

\[
Var(y)=\sigma_A^2+\sigma_D^2+\sigma_E^2
\]

\begin{verbatim}
model_ADE <- lme(
  bmi~1,
  random=list(
    famid=pdDiag(~Acoef-1),
    sibshipid=pdDiag(~Dcoef-1)
  ),
  data=longdata, method="REML")
\end{verbatim}

%====================================================
\section{Extracting Variance Components}

\begin{verbatim}
VarCorr(model_AE)
VarCorr(model_ACE)
VarCorr(model_ADE)
\end{verbatim}

%====================================================
\section{Model Comparison}

\begin{verbatim}
anova(model_AE, model_ACE)
anova(model_AE, model_ADE)
\end{verbatim}

%====================================================
\section{Derived Quantities}

\[
h^2=\frac{V_A}{V_P},\quad
c^2=\frac{V_C}{V_P},\quad
d^2=\frac{V_D}{V_P},\quad
e^2=\frac{V_E}{V_P}
\]

%====================================================
\section{Conclusion}

\begin{itemize}
\item Trio data estimate additive genetic variance and one familial component.
\item Siblings enable separation of shared environment and dominance.
\item Linear mixed models provide a practical implementation in R.
\end{itemize}

\end{document}
